% Part: computability
% Chapter: recursive-functions
% Section: pr-functions

\documentclass[../../../include/open-logic-section]{subfiles}

\begin{document}

\olfileid{cmp}{rec}{prf}
\olsection{આદિમ પુનરાવર્તી વિધેયો}

આદિમ પુનરાવર્તન અને સંયોજન વડે પહેલેથી મળેલા વિધેયોમાંથી નવાં વિધેયો
કેવી રીતે વ્યાખ્યાયિત થાય છે, તે ફરી નોંધીએ.

\begin{defn}
  \ollabel{defn:primitive-recursion} ધારો કે $f$ એ $k$-આર્ગ્યુમેન્ટવાળું
  વિધેય છે ($k\ge 1$) અને $g$ એ $(k+2)$-આર્ગ્યુમેન્ટવાળું વિધેય છે.
  $f$ અને $g$માંથી \emph{આદિમ પુનરાવર્તન દ્વારા વ્યાખ્યાયિત વિધેય}
  એટલે નીચેના સમીકરણોથી વ્યાખ્યાયિત $(k+1)$-આર્ગ્યુમેન્ટવાળું
  વિધેય~$h$:
  \begin{align*}
    h(x_0,\dots,x_{k-1},0) & =  f(x_0,\dots,x_{k-1}) \\
    h(x_0,\dots,x_{k-1},y+1) & = g(x_0,\dots,x_{k-1}, y, h(x_0,\dots,x_{k-1}, y))
  \end{align*}
\end{defn}

\begin{defn}
  \ollabel{defn:composition} ધારો કે $f$ એ $k$-આર્ગ્યુમેન્ટવાળું
  વિધેય છે, અને $g_0$, \dots, $g_{k-1}$ એ બધાં
  $n$-આર્ગ્યુમેન્ટવાળાં એવાં $k$ વિધેયો છે. $f$ તથા
  $g_0$, \dots,~$g_{k-1}$માંથી \emph{સંયોજન દ્વારા વ્યાખ્યાયિત વિધેય}
  એટલે નીચેના સમીકરણથી વ્યાખ્યાયિત $n$-આર્ગ્યુમેન્ટવાળું વિધેય~$h$:
  \[
  h(x_0, \dots, x_{n-1}) =
  f(g_0(x_0, \dots, x_{n-1}), \dots, g_{k-1}(x_0, \dots, x_{n-1})).
  \]
\end{defn}

$\Succ$ અને દરેક પ્રાકૃતિક સંખ્યા $n$ તથા $i < n$ માટેનાં પ્રક્ષેપ
વિધેયો
\[
\Proj{n}{i}(x_0,\dots,x_{n-1}) = x_i,
\]
ઉપરાંત, આદિમ પુનરાવર્તી વિધેયોમાં આપણે $\Zero(x) = 0$ વિધેયનો પણ
સમાવેશ કરીશું.

\begin{defn}
  આદિમ પુનરાવર્તી વિધેયોનો ગણ એટલે $\Nat^n$માંથી $\Nat$માં જતાં
  વિધેયોનો એવો ગણ, જેને નીચેની શરતોથી આગમનાત્મક રીતે વ્યાખ્યાયિત
  કરીએ છીએ:
  \begin{enumerate}
  \item $\Zero$ આદિમ પુનરાવર્તી છે.
  \item $\Succ$ આદિમ પુનરાવર્તી છે.
  \item દરેક પ્રક્ષેપ વિધેય $\Proj{n}{i}$ આદિમ પુનરાવર્તી છે.
  \item જો $f$ એ $k$-આર્ગ્યુમેન્ટવાળું આદિમ પુનરાવર્તી વિધેય હોય અને
    $g_0$, \dots,~$g_{k-1}$ એ $n$-આર્ગ્યુમેન્ટવાળાં આદિમ
    પુનરાવર્તી વિધેયો હોય, તો $g_0$, \dots,~$g_{k-1}$ સાથે $f$નું
    સંયોજન પણ આદિમ પુનરાવર્તી છે.
  \item જો $f$ એ $k$-આર્ગ્યુમેન્ટવાળું આદિમ પુનરાવર્તી વિધેય હોય અને
    $g$ એ $k+2$-આર્ગ્યુમેન્ટવાળું આદિમ પુનરાવર્તી વિધેય હોય, તો
    $f$ અને $g$માંથી આદિમ પુનરાવર્તન દ્વારા વ્યાખ્યાયિત વિધેય પણ
    આદિમ પુનરાવર્તી છે.
\end{enumerate}
\end{defn}

\begin{explain}
વધુ ટૂંકમાં, આદિમ પુનરાવર્તી વિધેયોનો ગણ એવો સૌથી નાનો ગણ છે જેમાં
$\Zero$, $\Succ$ અને પ્રક્ષેપ વિધેયો~$\Proj{n}{j}$ છે અને જે
સંયોજન તથા આદિમ પુનરાવર્તન હેઠળ બંધ છે.

આદિમ પુનરાવર્તી વિધેયોનો ગણ ``તબક્કાઓ'' વડે પણ વર્ણવી શકાય. $S_0$
એ શરૂઆતનાં વિધેયોનો ગણ દર્શાવે: $\Zero$, $\Succ$ અને પ્રક્ષેપ
વિધેયો. આ તબક્કા~$0$નાં આદિમ પુનરાવર્તી વિધેયો છે. તબક્કો $S_i$
વ્યાખ્યાયિત થયા પછી, $S_i$માંનાં વિધેયો પર સંયોજન અથવા આદિમ
પુનરાવર્તનનો એક વાર ઉપયોગ કરીને મળતાં બધાં વિધેયોનો ગણ $S_{i+1}$
માનીએ. ત્યારે
\[
S = \bigcup_{i \in \Nat} S_i
\]
એ બધાં આદિમ પુનરાવર્તી વિધેયોનો ગણ છે.
\end{explain}

હવે તપાસીએ કે $\Add$ આદિમ પુનરાવર્તી વિધેય છે.

\begin{prop}
  સરવાળાનું વિધેય $\Add(x,y) = x+y$ આદિમ પુનરાવર્તી છે.
\end{prop}

\begin{proof}
આપણે $f$ અને~$g$ એવાં બે વિધેયોના આધારે $\Add$ની આદિમ પુનરાવર્તી
વ્યાખ્યા પહેલેથી આપી છે, જે \olref{defn:primitive-recursion}ના
માળખાને અનુસરે છે:
\begin{align*}
  \Add(x_0, 0) & = f(x_0) = x_0\\
  \Add(x_0, y+1) & = g(x_0, y, \Add(x_0, y)) =
  \Succ(\Add(x_0, y))
\end{align*}
તેથી $f$ અને $g$ પણ આદિમ પુનરાવર્તી હોય તો $\Add$ આદિમ
પુનરાવર્તી છે. $f(x_0) = x_0 = \Proj{1}{0}(x_0)$ છે અને પ્રક્ષેપ
વિધેયોને આદિમ પુનરાવર્તી ગણીએ છીએ, તેથી $f$ આદિમ પુનરાવર્તી છે.
$g$ એ $g(x_0, y, z)$ એવાં ત્રણ આર્ગ્યુમેન્ટવાળું વિધેય છે, જે આ
રીતે વ્યાખ્યાયિત છે:
\[
g(x_0, y, z) = \Succ(z).
\]
આટલાથી હજી $g$ આદિમ પુનરાવર્તી છે એવું સાબિત થતું નથી, કારણ કે
$g$ અને $\Succ$ એકસરખાં વિધેયો નથી: $\Succ$ને એક આર્ગ્યુમેન્ટ છે,
જ્યારે $g$ને ત્રણ આર્ગ્યુમેન્ટ હોવા જોઈએ. પરંતુ સંયોજન દ્વારા
$g$ને ઔપચારિક રીતે આમ વ્યાખ્યાયિત કરી શકીએ:
\[
g(x_0, y, z) = \Succ(\Proj{3}{2}(x_0, y, z))
\]
$\Succ$ અને $\Proj{3}{2}$ને આદિમ પુનરાવર્તી વિધેયો ગણીએ છીએ;
તેમના સંયોજનથી આ વિધેય વ્યાખ્યાયિત થાય છે, તેથી $g$ પણ આદિમ
પુનરાવર્તી છે.
\end{proof}

\begin{prop}
  \ollabel{prop:mult-pr}
  ગુણાકારનું વિધેય $\Mult(x,y) = x \cdot y$ આદિમ પુનરાવર્તી છે.
\end{prop}

\begin{proof}
  અભ્યાસ માટે.
\end{proof}

\begin{prob}
  $\Mult$ની આદિમ પુનરાવર્તી વ્યાખ્યા
  \olref[cmp][rec][prf]{defn:primitive-recursion}માં જરૂરી માળખામાં
મૂકી શકાય છે અને સંબંધિત વિધેયો $f$ તથા~$g$ આદિમ પુનરાવર્તી છે,
એવું બતાવીને \olref[cmp][rec][prf]{prop:mult-pr} સાબિત કરો.
\end{prob}

\begin{ex}
આદિમ પુનરાવર્તી વ્યાખ્યાનું આપણું સૌપ્રથમ ઉદાહરણ આ હતું:
\begin{align*}
h(0) & =  1 \\
h(y+1) & =  2 \cdot h(y).
\end{align*}
આ વિધેય \olref{defn:primitive-recursion}માં જરૂરી માળખામાં બંધબેસતું
નથી, કારણ કે અહીં $k=0$ છે. વ્યાખ્યામાં અચળ મૂલ્યો $1$ અને $2$ પણ
આવે છે. પહેલી મુશ્કેલી દૂર કરવા એક બિનવપરાયેલું આર્ગ્યુમેન્ટ ઉમેરીને
વિધેય~$h'$ વ્યાખ્યાયિત કરીએ:
\begin{align*}
h'(x_0, 0) & =  f(x_0) = 1 \\
h'(x_0, y+1) & =  g(x_0, y, h'(x_0, y)) = 2 \cdot h'(x_0, y).
\end{align*}
સંયોજન વડે $\Succ$ અને $\Zero$માંથી $f(x_0) = 1$ વ્યાખ્યાયિત
કરી શકાય: $f(x_0) = \Succ(\Zero(x_0))$. $g'(z) = 2 \cdot z$
અને પ્રક્ષેપ વિધેયોમાંથી સંયોજન વડે $g$ વ્યાખ્યાયિત કરી શકાય:
\begin{align*}
g(x_0, y, z) & = g'(\Proj{3}{2}(x_0, y, z))
\intertext{અને $g'$ને પણ સંયોજન વડે આમ વ્યાખ્યાયિત કરી શકાય:}
g'(z) & = \Mult(g''(z), \Proj{1}{0}(z))
\intertext{તથા}
g''(z) & = \Succ(f(z)),
\end{align*}
અહીં $f$ ઉપર મુજબ છે: $f(z) = \Succ(\Zero(z))$. હવે~$h'$
મળી ગયું છે, એટલે ફરી સંયોજન વાપરીને $h(y) =
h'(\Proj{1}{0}(y),\Proj{1}{0}(y))$ રાખી શકીએ. આથી $h$ને
મૂળભૂત વિધેયોમાંથી સંયોજનો અને આદિમ પુનરાવર્તનોની શ્રેણી વડે
વ્યાખ્યાયિત કરી શકાય છે; એટલે $h$ આદિમ પુનરાવર્તી છે.
\end{ex}

\end{document}
